\documentclass[10pt,a4paper]{amsart}
\usepackage[margin=19mm]{geometry}
\usepackage{amsmath,amssymb,mathtools}
\usepackage[scr=rsfs,cal=euler]{mathalfa}
\usepackage{xfrac}
\usepackage[unicode,hidelinks,bookmarks=false]{hyperref}
\newcommand{\ie}{i.e.\ }
\newcommand{\cf}{cf.\ }
\newcommand{\resp}{respectively\ }
\newcommand{\Subsection}{\subsection}
\providecommand{\nobreakdash}{\nobreak\hbox{-}}
\setlength{\emergencystretch}{2em}
\multlinegap0pt
\usepackage{bm}
\usepackage{bbm}
\usepackage{dsfont}
\usepackage{xspace}
\usepackage{cancel}
\usepackage{mathtools}
\usepackage{amsthm}
\makeatletter
\@ifundefined{theorem}{\newtheorem{theorem}{Theorem}}{}
\@ifundefined{definition}{\newtheorem{definition}[theorem]{Definition}}{}
\makeatother
\title[The minimal periodicity for integral bases of pure number fi]{The minimal periodicity for integral bases of pure number fields}
\author{Khai-Hoan Nguyen-Dang}
\date{Source: arXiv:2509.09457v3 (CC BY 4.0)}
\begin{document}
\maketitle

\noindent\emph{Source and licence.} The minimal periodicity for integral bases of pure number fields. Version arXiv:2509.09457v3 (CC BY 4.0), distributed under the Creative Commons Attribution 4.0 International licence, as recorded in the arXiv licence metadata for this exact version and shown on the abstract page https://arxiv.org/abs/2509.09457v3. Full rights notice: licenses/arxiv-2509.09457-CC-BY-4.0.txt.\par\smallskip

\noindent\textit{Editorial scope.} Counted unit: the Theorem whose source label is \texttt{thm:equidistribution} in the article (source0.tex lines 845--884, source label \texttt{thm:equidistribution}), delivered together with the article's own complete original proof of it (source0.tex lines 886--1044, one \texttt{proof} environment of 159 lines). No mathematics is written, repaired or re-derived: statement and proof are the article's own text line for line. Required context retained uncounted: def:admissibilities (lines 773--795); eq:mobius-nfree (lines 825--827). Every definition, display and label of those lines is retained unchanged. Omitted from this package and not redistributed: 1--772, 796--824, 828--844, 885--885, 1045--1241. Nothing inside a retained excerpt is omitted or altered: every retained body is delivered exactly as the article's lines are written, so no omission receipt of any kind accompanies this package. Citations: the retained excerpts carry no citation command at all, so no bibliography entry of the article is transcribed and no reference is redistributed. Reference adaptation: none was needed and none was made; every reference inside the delivered unit resolves inside it, so no unresolved reference or citation is produced. Printed slips kept literal: no printed word, symbol, hypothesis or number of the counted statement or of its proof was altered. Layout and class: the article's own class is \texttt{amsart} with packages including amsmath, amsthm, amssymb, amsfonts, mathtools, geometry, microtype, hyperref; the wrapper is the lane's pinned \texttt{amsart} editorial wrapper at 10pt on A4, which restates the theorem-like environments the retained text needs and adds the article's own macro definitions that the retained text needs (none), with extra wrapper packages (bm, bbm, dsfont, xspace, cancel, mathtools). The delivered numbering is the unit's own. Assets: no retained span contains a figure environment, a graphics-inclusion command, a PDF-page inclusion, a table or a TikZ picture; no image file is copied into this package, the package declares no asset dependency and no image file is redistributed with it. Licence: version arXiv:2509.09457v3 of this article is distributed under the Creative Commons Attribution 4.0 International licence, as recorded in the arXiv licence metadata for this exact version and shown on the abstract page https://arxiv.org/abs/2509.09457v3. A full rights notice is delivered at licenses/arxiv-2509.09457-CC-BY-4.0.txt. Characters: every retained excerpt is pure ASCII, so the delivered engine, XeLaTeX with its default Latin Modern text font, reports no missing character.
\begin{definition}[Admissibility notion]\label{def:admissibilities}
Fix $q\ge1$ and $n\ge2$, and let $r$ be a residue class modulo $q$.  If
$p^\alpha\parallel q$, define the $p$--adic valuation of $r$ modulo $q$ by
\[
 v_{p,q}(r):=\max\{\,0\le s\le \alpha:\ r\equiv0\pmod{p^s}\,\}.
\]
Equivalently, if $\widetilde r$ is any integer representative of the class $r$, then
\[
 v_{p,q}(r)=\min\{v_p(\widetilde r),\alpha\},
\]
with the convention $v_p(0)=+\infty$.  Thus $v_{p,q}(r)$ is independent of the chosen
representative.  When $q$ is fixed, we write simply $v_p(r)$ for $v_{p,q}(r)$.
Moreover, if $g\mid q$, then the notation $g\mid r$ means $r\equiv0\pmod g$; this is again
well defined for a residue class modulo $q$.
\begin{itemize}
\item[(a)] We say that $r$ is \emph{$n$--admissible} (weak form) if for every
$p^\alpha\parallel q$ with $\alpha\ge n$ we have $v_p(r)<n$.  Equivalently, the progression
$a\equiv r\pmod q$ does not force the divisibility $p^n\mid a$ for any prime $p$.
\item[(b)] We say that $r$ is \emph{strictly $n$--admissible} if for every $p^\alpha\parallel q$
we have $v_p(r)<\min\{\alpha,n\}$.  This is stronger than (a) and is exactly the condition
that makes the main term below independent of $r$.
\end{itemize}
\end{definition}

\begin{equation}\label{eq:mobius-nfree}
\mathbf 1_{\text{$n$--free}}(a)=\sum_{d^n\mid a}\mu(d).
\end{equation}

\begin{theorem}[Equidistribution with explicit $r$-dependence]\label{thm:equidistribution}
Fix $n\ge2$. For every $q\ge1$, every residue class $r\pmod q$, and every real $X\ge1$,
\begin{equation}\label{eq:Fn-mainterm}
\begin{aligned}
\mathcal F_n(X;q,r)
={}&
\frac{2X}{q}\cdot\frac{1}{\zeta(n)}
\prod_{p\mid q}\Bigl(1-\frac1{p^{\,n}}\Bigr)^{-1}\\
&\times
\prod_{p^\alpha\parallel q}
\Bigl(1-\mathbf 1_{\{v_p(r)\ge \min(\alpha,n)\}}
      p^{\min(\alpha,n)-n}\Bigr)
+E_n(X;q,r),
\end{aligned}
\end{equation}
where
\begin{equation}\label{eq:Fn-error-explicit}
\bigl|E_n(X;q,r)\bigr|
\le
\left(2+\frac{2n}{n-1}\right)X^{1/n}.
\end{equation}
In particular, $E_n(X;q,r)=O_n(X^{1/n})$, uniformly in $q$ and $r$.
For fixed $q$ and $r$, formula \eqref{eq:Fn-mainterm} is therefore an asymptotic formula as
$X\to\infty$, because $X^{1/n}=o(X)$.

Here $v_p(r)=v_{p,q}(r)$ is the residue-class valuation from Definition~\ref{def:admissibilities}.
In particular:
\begin{itemize}
\item If $r$ is \emph{not} $n$--admissible, then $\mathcal F_n(X;q,r)=0$ for every $X$.
\item If $r$ is \emph{strictly $n$--admissible}, then the $r$-dependent product in
\eqref{eq:Fn-mainterm} equals $1$, and
\[
\mathcal F_n(X;q,r)=
\frac{2X}{q}\cdot\frac{1}{\zeta(n)}
\prod_{p\mid q}\Bigl(1-\frac1{p^{\,n}}\Bigr)^{-1}
+O_n\!\bigl(X^{1/n}\bigr),
\]
independently of $r$.
\end{itemize}
\end{theorem}

\begin{proof}
Put $Y=X^{1/n}$. By definition and by \eqref{eq:mobius-nfree},
\[
\mathcal F_n(X;q,r)
=\sum_{\substack{0<|a|\le X\\ a\equiv r\ (\mathrm{mod}\ q)}}
 \sum_{d^n\mid a}\mu(d).
\]
If $d^n\mid a$ and $0<|a|\le X$, then $d\le Y$. Hence the double sum is finite, and we may
interchange the order of summation:
\begin{equation}\label{eq:Fn-swap}
\mathcal F_n(X;q,r)
=\sum_{1\le d\le Y}\mu(d)\,
\#\bigl\{\,0<|a|\le X:\ a\equiv r\pmod q,\ d^n\mid a\,\bigr\}.
\end{equation}

Fix such a $d$ and put
\[
g_d:=\gcd(q,d^n),\qquad
L_d:=\operatorname{lcm}(q,d^n)=\frac{qd^n}{g_d}.
\]
The simultaneous congruences
\[
a\equiv r\pmod q,
\qquad a\equiv0\pmod{d^n}
\]
are compatible if and only if $r\equiv0\pmod{g_d}$, i.e. if and only if $g_d\mid r$ in the
residue-class sense of Definition~\ref{def:admissibilities}. If they are compatible, their
solutions form one residue class modulo $L_d$.

We now make the counting error completely explicit. For any modulus $L\ge1$ and any residue
class $c\pmod L$, choose an integer representative $c_0$. The relevant count is
\[
\left\lfloor\frac{X-c_0}{L}\right\rfloor
-\left\lceil\frac{-X-c_0}{L}\right\rceil+1,
\]
and therefore differs from $2X/L$ by at most $1$. Removing the integer $0$ changes this count
by at most one. Consequently, for every $d$ there is a real number $\varepsilon_d(X;q,r)$ with
$|\varepsilon_d(X;q,r)|\le2$ such that
\begin{equation}\label{eq:single-progression-count}
\begin{aligned}
&\#\bigl\{\,0<|a|\le X:\ a\equiv r\pmod q,\ d^n\mid a\,\bigr\}\\
&\qquad=
\mathbf 1_{\{g_d\mid r\}}
\left(\frac{2X}{L_d}+\varepsilon_d(X;q,r)\right)\\
&\qquad=
\mathbf 1_{\{g_d\mid r\}}
\left(\frac{2X}{q}\frac{g_d}{d^n}+\varepsilon_d(X;q,r)\right).
\end{aligned}
\end{equation}
The second equality in \eqref{eq:single-progression-count} is exact, since
$L_d=qd^n/g_d$; no term is absorbed into the error.

Substituting \eqref{eq:single-progression-count} into \eqref{eq:Fn-swap} gives
\begin{equation}\label{eq:Fn-SX}
\mathcal F_n(X;q,r)=\frac{2X}{q}S_X(q,r)+E_1(X;q,r),
\end{equation}
where
\[
S_X(q,r):=
\sum_{1\le d\le Y}
\mu(d)\frac{\gcd(d^n,q)}{d^n}\,
\mathbf 1_{\{\gcd(d^n,q)\mid r\}}
\]
and
\[
E_1(X;q,r):=
\sum_{1\le d\le Y}
\mu(d)\mathbf 1_{\{g_d\mid r\}}\varepsilon_d(X;q,r).
\]
Since $|\mu(d)|\le1$ and $|\varepsilon_d(X;q,r)|\le2$,
\begin{equation}\label{eq:E1-bound}
|E_1(X;q,r)|\le2\lfloor Y\rfloor\le2Y.
\end{equation}
The constant $2$ in this estimate is absolute; in particular, it is independent of $q$ and $r$.

Consider now the absolutely convergent series
\[
S(q,r):=
\sum_{d\ge1}
\mu(d)\frac{\gcd(d^n,q)}{d^n}\,
\mathbf 1_{\{\gcd(d^n,q)\mid r\}}.
\]
Indeed, $\gcd(d^n,q)\le q$ and $n\ge2$. Moreover,
\[
|S(q,r)-S_X(q,r)|
\le q\sum_{d>Y}\frac1{d^n}.
\]
Since $Y\ge1$ and $t\mapsto t^{-n}$ is decreasing,
\[
\sum_{d>Y}\frac1{d^n}
\le Y^{-n}+\int_Y^\infty t^{-n}\,dt
\le \frac{n}{n-1}Y^{1-n}.
\]
Therefore
\begin{equation}\label{eq:tail-bound}
\left|\frac{2X}{q}\bigl(S(q,r)-S_X(q,r)\bigr)\right|
\le\frac{2n}{n-1}XY^{1-n}
=\frac{2n}{n-1}X^{1/n}.
\end{equation}
Combining \eqref{eq:Fn-SX}, \eqref{eq:E1-bound}, and \eqref{eq:tail-bound}, we obtain
\begin{equation}\label{eq:Fn-S}
\mathcal F_n(X;q,r)=\frac{2X}{q}S(q,r)+E_n(X;q,r),
\end{equation}
with the explicit bound \eqref{eq:Fn-error-explicit}. This also explains why the error depends
only on $n$: the factor $q$ in the tail estimate cancels against the prefactor $1/q$.

It remains to compute $S(q,r)$. The summand is multiplicative in $d$. Indeed, if
$\gcd(d_1,d_2)=1$, then
\[
\gcd((d_1d_2)^n,q)=\gcd(d_1^n,q)\gcd(d_2^n,q),
\]
the two factors on the right are coprime, and hence
\[
\mathbf 1_{\{\gcd((d_1d_2)^n,q)\mid r\}}
=
\mathbf 1_{\{\gcd(d_1^n,q)\mid r\}}
\mathbf 1_{\{\gcd(d_2^n,q)\mid r\}}.
\]
Absolute convergence therefore permits an Euler-product factorisation. Since $\mu(d)$ is
supported on squarefree integers, each prime contributes only the choices $p\nmid d$ and
$p\mid d$.

If $p^\alpha\parallel q$, the local term corresponding to $p\mid d$ is
\[
\mu(p)\,p^{\min(n,\alpha)-n}\,
\mathbf 1_{\{v_p(r)\ge\min(n,\alpha)\}},
\]
so the local factor is
\[
1-
\mathbf 1_{\{v_p(r)\ge\min(n,\alpha)\}}
 p^{\min(n,\alpha)-n}.
\]
If $p\nmid q$, the compatibility condition is automatic and the local factor is
$1+\mu(p)p^{-n}=1-p^{-n}$. Hence
\[
S(q,r)=
\prod_{p\nmid q}(1-p^{-n})
\prod_{p^\alpha\parallel q}
\Bigl(1-
\mathbf 1_{\{v_p(r)\ge\min(\alpha,n)\}}
 p^{\min(\alpha,n)-n}\Bigr).
\]
Finally,
\[
\prod_{p\nmid q}(1-p^{-n})
=
\prod_p(1-p^{-n})\prod_{p\mid q}(1-p^{-n})^{-1}
=
\frac1{\zeta(n)}\prod_{p\mid q}(1-p^{-n})^{-1}.
\]
Substitution into \eqref{eq:Fn-S} proves \eqref{eq:Fn-mainterm}.

If $r$ is not $n$--admissible, then for some $p^\alpha\parallel q$ with $\alpha\ge n$ one has
$v_p(r)\ge n$. Thus every integer $a\equiv r\pmod q$ is divisible by $p^n$, and therefore
$\mathcal F_n(X;q,r)=0$. If $r$ is strictly $n$--admissible, then
$v_p(r)<\min(\alpha,n)$ for every $p^\alpha\parallel q$, so all indicators in the second
product vanish. This proves the stated $r$-independent main term.
\end{proof}

\end{document}
